\documentclass[12pt,a4paper]{article}
\usepackage[utf8]{inputenc}
\usepackage[T1]{fontenc}
\usepackage[spanish,es-noshorthands]{babel}
\usepackage{amsmath,amssymb,geometry,fancyhdr,xcolor,listings,tikz,array}
\geometry{margin=2.4cm,headheight=16pt}
\pagestyle{fancy}\fancyhf{}\lhead{UPCT}\rhead{Cálculo diferencial}\lfoot{Prácticas de wxMaxima}\rfoot{\thepage}
\definecolor{azul}{RGB}{23,65,105}
\lstset{basicstyle=\ttfamily\small,backgroundcolor=\color{black!4},frame=single,rulecolor=\color{black!15},breaklines=true,columns=fullflexible}

\begin{document}
\begin{center}{\large Práctica de laboratorio IV}\\[5pt]{\LARGE\bfseries\color{azul} Extremos condicionados}\end{center}
\noindent Nombre y apellidos: \hrulefill\par\medskip
\noindent\textbf{Objetivos.} Obtener y clasificar extremos bajo restricciones de igualdad, distinguir condiciones necesarias y suficientes y resolver problemas geométricos con wxMaxima.
\section{Extremos sobre un conjunto de restricciones}
Sea $f:U\subset\mathbb R^n\to\mathbb R$, con $U$ abierto, y sea $M=\{X\in U:g_1(X)=\cdots=g_m(X)=0\}$. Un mínimo relativo condicionado en $P\in M$ satisface $f(P)\le f(X)$ para todo $X\in M$ suficientemente próximo a $P$. Para un máximo se invierte la desigualdad. El extremo es estricto si la desigualdad es estricta para $X\ne P$.

Si $M$ es no vacío y compacto y $f$ es continua, el teorema de Weierstrass garantiza máximos y mínimos absolutos. No garantiza que todos los candidatos sean extremos.
\section{Condición necesaria: multiplicadores de Lagrange}
Supongamos $f,g_i\in C^1(U)$ y que los gradientes de las restricciones son linealmente independientes en $P$. El espacio tangente es $T_PM=\ker Dg(P)$. En un extremo, $Df(P)v=0$ para todo $v\in T_PM$. Por tanto,
\[\boxed{\nabla f(P)=\sum_{i=1}^m\lambda_i\nabla g_i(P),\quad g(P)=0.}\]
Introducimos $\mathcal L(X,\lambda)=f(X)-\sum_i\lambda_i g_i(X)$. En dos variables y una restricción, se resuelve $\mathcal L_x=\mathcal L_y=g=0$. La condición es necesaria, no suficiente.

\textbf{Ejemplo guiado.} Optimizar $f=x+2y$ sobre $x^2+y^2=1$. Los candidatos son $\pm(1,2)/\sqrt5$; sus valores son $\pm\sqrt5$. Comparar los valores y justificar la existencia mediante compacidad.
\begin{lstlisting}
kill(all)$
f:x+2*y$ g:x^2+y^2-1$
L:f-lam*g$
S:solve([diff(L,x)=0,diff(L,y)=0,g=0],[x,y,lam]);
map(lambda([s],radcan(subst(s,f))),S);
wxplot2d([parametric,cos(t),sin(t),[t,0,2*%pi]],[x,-1.2,1.2],[y,-1.2,1.2])$
\end{lstlisting}
\section{Condiciones de segundo orden}
Sean $f,g_i\in C^2(U)$ y $P$ un punto regular que satisface las ecuaciones de Lagrange. Se calcula la hessiana $H_X\mathcal L$ respecto de las variables originales, manteniendo los multiplicadores fijos. La forma cuadrática relevante es
\[Q(v)=v^TH_X\mathcal L(P,\lambda)v,\qquad v\in T_PM.\]
Para un mínimo (máximo) es necesario que $Q$ sea semidefinida positiva (negativa) sobre el espacio tangente. Si es definida positiva (negativa) sobre dicho espacio, hay mínimo (máximo) relativo estricto. Si toma ambos signos, no hay extremo; si es semidefinida singular, el criterio no decide.

En dos variables, con $\nabla g\ne0$, basta evaluar $Q$ en $v=(-g_y,g_x)$. Si $Q(v)>0$ hay mínimo estricto; si $Q(v)<0$ hay máximo estricto. Si vale cero, se requieren otros métodos.

La justificación se obtiene derivando dos veces $f(\gamma(t))$ para una curva admisible $\gamma$ con $\gamma(0)=P$: al usar $g(\gamma(t))=0$ y la ecuación de Lagrange, resulta $(f\circ\gamma)''(0)=Q(\gamma'(0))$. Una carta local y Taylor justifican la suficiencia.

\textbf{Hessiana orlada.} En dos variables,
\[B=\begin{pmatrix}0&g_x&g_y\\g_x&\mathcal L_{xx}&\mathcal L_{xy}\\g_y&\mathcal L_{yx}&\mathcal L_{yy}\end{pmatrix},\qquad \det B=-Q(v).\]
\textbf{Ejemplo.} $f=x^2+2y^2$ sobre la circunferencia unidad. Clasificar los cuatro candidatos mediante $Q$.
\begin{lstlisting}
f:x^2+2*y^2$ g:x^2+y^2-1$ L:f-lam*g$
S:solve([diff(L,x)=0,diff(L,y)=0,g=0],[x,y,lam]);
H:matrix([diff(L,x,2),diff(diff(L,x),y)],[diff(diff(L,y),x),diff(L,y,2)])$
v:matrix([-diff(g,y)],[diff(g,x)])$
q:expand((transpose(v).H.v)[1,1]);
for s in S do print(s,subst(s,f),subst(s,q))$
wxplot2d(cos(t)^2+2*sin(t)^2,[t,0,2*%pi])$
\end{lstlisting}
\section{Parametrización y casos singulares}
La parametrización de la restricción proporciona un método alternativo. Por ejemplo, si $g=y-x^2=0$, se estudia $F(t)=f(t,t^2)$. Para $f=y-x^2+x^4$, resulta $F(t)=t^4$: el origen es mínimo estricto aunque la segunda derivada restringida sea nula.

Si los gradientes de las restricciones no son independientes, el teorema de Lagrange no se aplica. Por ejemplo, $g=y^2-x^3=0$ y $f=x$: el origen es mínimo absoluto estricto, pero $\nabla g(0,0)=0$ y $\nabla f(0,0)=(1,0)$. La parametrización $(t^2,t^3)$ demuestra el resultado. Los puntos singulares deben investigarse por separado.
\begin{lstlisting}
f:y-x^2+x^4$ F:expand(subst([x=t,y=t^2],f));
diff(F,t,2);
f:x$ g:y^2-x^3$
subst([x=0,y=0],[diff(g,x),diff(g,y)]);
subst([x=t^2,y=t^3],[f,g]);
\end{lstlisting}
\section{Varias restricciones y aplicaciones}
Para $f=x^2+y^2+z^2$ sujeto a $x+y+z=3$ y $x-y=0$, ambos gradientes son independientes. Las ecuaciones de Lagrange dan $(1,1,1)$, con valor $3$. El espacio tangente está generado por $(1,1,-2)$; la hessiana de $\mathcal L$ es $2I$, y el mínimo es estricto. También es absoluto, pues sobre la recta admisible $f(t,t,3-2t)=3+6(t-1)^2$.

\textbf{Aplicación geométrica.} Un rectángulo de lados positivos $x,y$ tiene perímetro $20$. Maximizar $xy$ sujeto a $x+y=10$. El cuadrado de lado $5$ da el máximo absoluto $25$, como se comprueba por sustitución. El dominio de lados positivos no es compacto; no hay mínimo absoluto, aunque el ínfimo es cero.
\begin{lstlisting}
f:x^2+y^2+z^2$ g1:x+y+z-3$ g2:x-y$
L:f-lam*g1-mu*g2$
solve([diff(L,x)=0,diff(L,y)=0,diff(L,z)=0,g1=0,g2=0],[x,y,z,lam,mu]);
expand(subst([x=t,y=t,z=3-2*t],f));
f:x*y$ g:x+y-10$ L:f-lam*g$
solve([diff(L,x)=0,diff(L,y)=0,g=0],[x,y,lam]);
wxplot2d(x*(10-x),[x,0,10])$
\end{lstlisting}
\section{Trabajo individual}
En cada ejercicio, determinar el conjunto admisible, comprobar regularidad, obtener candidatos, clasificar y justificar si son absolutos. Si el criterio no decide, utilizar una parametrización o una desigualdad.
\begin{enumerate}
\item Optimizar $xy$ sobre $x^2+y^2=4$.
\item Optimizar $x^2+4y^2$ sobre $x+y=5$.
\item Optimizar $x$ sobre $x^2+4y^2=4$.
\item Estudiar $x^4$ sobre $y=0$.
\item Estudiar $y$ sobre $x^2+y^2=0$.
\item Hallar el punto del plano $x+2y+2z=9$ más próximo al origen.
\end{enumerate}
\section{Cuestiones de comprensión}
\begin{enumerate}
\item ¿Por qué no basta con anular el gradiente de la función objetivo?
\item ¿Por qué se utiliza la hessiana de la lagrangiana y no la de la función objetivo?
\item ¿Qué diferencia existe entre un punto singular de la restricción y un punto estacionario degenerado?
\item ¿Es suficiente comparar los valores de los candidatos si el conjunto no es compacto?
\end{enumerate}
\end{document}
